% Eigenständige Lesefassung der vollständigen Ganzzahlkonstruktion.
% Formale Ableitungen stehen im zugehörigen Ergänzungsband.
\par\Needspace{10\baselineskip}
\hypertarget{integers.reading.problem}{}
\section*{1. Warum die natürlichen Zahlen nicht ausreichen}

Die Gleichung
\[
  x+5=2
\]
hat in den natürlichen Zahlen keine Lösung. Wer bereits mit ganzen Zahlen
rechnet, schreibt dafür \(x=-3\). In dieser Lesefassung soll erklärt
werden, wie ein Zahlenbereich entsteht, in dem diese Antwort einen
bestimmten Sinn hat. Ein neues Zahlzeichen allein genügt dafür nicht.
Wir benötigen Objekte, eine Gleichheit zwischen ihnen und Rechenregeln,
die mit dieser Gleichheit verträglich sind.

Das Ziel geht über die Möglichkeit des Subtrahierens hinaus. Wir wollen
ganze Zahlen addieren und multiplizieren, ihre Größe vergleichen und
mit Induktion über sie argumentieren können. Am Ende sollen alle
Eigenschaften begründet sein, die Band~17 anschließend als Axiome der
ganzen Zahlen verwendet. Erst dann dürfen wir die konkrete Konstruktion
verlassen und nur noch mit diesen Axiomen arbeiten.

Unser Ausgangspunkt sind die natürlichen Zahlen einschließlich der Null.
Ihre Addition und Multiplikation, ihre Ordnung und ihre Induktion werden
vorausgesetzt. Dazu gehören insbesondere die Kürzungsgesetze für
natürliche Zahlen \(r,s,t\):
\[
  r+t=s+t\ \Longrightarrow\ r=s,
  \qquad
  r\leq s\ \Longleftrightarrow\ r+t\leq s+t.
\]
Auch die Beschreibung der natürlichen Ordnung durch einen Abstand ist
bereits bekannt:
\[
  r\leq s
  \quad\Longleftrightarrow\quad
  \text{es gibt ein }n\in\mathbb N\text{ mit }r+n=s.
\]
Wenn wir später Summanden kürzen oder einen solchen Abstand wählen,
verwenden wir diese Eigenschaften der natürlichen Zahlen. Eine noch
nicht konstruierte Subtraktion wird dabei nicht vorausgesetzt.

Die Konstruktion beginnt mit Paaren natürlicher Zahlen. Ein Paar
\((a,b)\) soll die Absicht ausdrücken, \(b\) von \(a\) abzuziehen.
Anschließend werden Paare, die dieselbe Differenz darstellen sollen, zu
Klassen zusammengefasst. Die Klassen werden die neuen Zahlen. Wir
begründen erst ihr Rechnen, danach ihre Ordnung und zuletzt die
Induktion in beide Richtungen.

Die \IntegersProofsLink{} enthalten die
vollständigen formalen Ableitungen bis zum Modellnachweis. Der vorliegende
Text erklärt dieselben mathematischen Schritte in zusammenhängender
Sprache. Verweise auf benannte Aussagen führen zur entsprechenden
Stelle in Band~17 oder seinem Ergänzungsband. Die spätere
\DifferencesReadingLink{} zu formalen Differenzen in Band~41
greift die additive Grundidee wieder auf und untersucht, unter welchen
allgemeinen Voraussetzungen sie funktioniert. Hier steht die gesamte
konkrete Zahlstruktur im Mittelpunkt.

\par\Needspace{10\baselineskip}
\hypertarget{integers.reading.classes}{}
\section*{2. Aus Paaren werden Zahlen}

Die Paare \((2,5)\), \((3,6)\) und \((0,3)\) sollen dieselbe
neue Zahl beschreiben. Die ersten beiden unterscheiden sich nur dadurch,
dass beide Einträge um eins erhöht wurden. Für eine gedachte Differenz
wäre das ohne Bedeutung. Um dies ohne Subtraktion auszudrücken,
vergleichen wir Kreuzsummen:
\[
  2+6=5+3,
  \qquad
  2+3=5+0.
\]
Das führt auf die Definition
\[
  (a,b)\sim(c,d)
  \quad\Longleftrightarrow\quad
  a+d=b+c
  \qquad(a,b,c,d\in\mathbb N).
\]
Die rechte Seite benutzt ausschließlich die natürliche Addition. Damit
ist die Relation bereits verständlich, bevor negative Zahlen existieren.
Formal steht sie in \FormulaRefAuto{IntPairRelDef}[def].

Wir wollen äquivalente Paare zu einem einzigen Objekt zusammenfassen.
Dafür muss \(\sim\) eine Äquivalenzrelation sein. Das bedeutet
Reflexivität, Symmetrie und Transitivität. Die ersten beiden Eigenschaften
folgen unmittelbar aus der Kommutativität der Addition: Es gilt
\(a+b=b+a\), also \((a,b)\sim(a,b)\); und aus
\(a+d=b+c\) folgt durch Umordnen und Umkehren
\(c+b=d+a\), also \((c,d)\sim(a,b)\).

\medskip\noindent\textbf{Was die Transitivität verlangt.}
Sind das erste und das zweite Paar äquivalent und ebenso das zweite
und das dritte, dann sollen auch das erste und das dritte äquivalent
sein. Für drei Paare \((a,b)\), \((c,d)\) und \((e,f)\) ist also
zu zeigen:
\[
  (a,b)\sim(c,d)\ \text{ und }\ (c,d)\sim(e,f)
  \quad\Longrightarrow\quad (a,b)\sim(e,f).
\]
Die beiden Voraussetzungen bedeuten nach unserer Definition
\[
  a+d=b+c,
  \qquad c+f=d+e.
\]
Die gewünschte Folgerung bedeutet \(a+f=b+e\). Um die beiden
Voraussetzungen miteinander zu verbinden, vergleichen wir zunächst
\((a+f)+d\) und \((b+e)+d\). Dabei verwenden wir nur die
Assoziativität und Kommutativität der natürlichen Addition sowie
die beiden vorausgesetzten Gleichungen:
\[
  \begin{aligned}
    (a+f)+d
      &=(a+d)+f\\
      &=(b+c)+f\\
      &=b+(c+f)\\
      &=b+(d+e)\\
      &=(b+e)+d.
  \end{aligned}
\]
Das natürliche Kürzungsgesetz entfernt den gemeinsamen Summanden \(d\).
Damit gilt \((a,b)\sim(e,f)\). Gerade in diesem Schritt trägt die
Kürzbarkeit der natürlichen Addition die Konstruktion. Die formale
Zusammenfassung ist \FormulaRefAuto{IntegerPairEquivalence}[thm].

\medskip\noindent\textbf{Eine Zahl ist eine Klasse von Darstellungen.}
Zu jedem Paar bilden wir seine Äquivalenzklasse
\[
  [a,b]:=\{(c,d)\in\mathbb N\times\mathbb N
                    \mid(c,d)\sim(a,b)\}.
\]
Das Paar \((a,b)\) heißt ein Repräsentant dieser Klasse. Es ist
ein Element der Klasse; die Klasse selbst enthält alle gleichwertigen
Darstellungen. Die Menge aller solchen Klassen ist der Quotient
\[
  \mathbb Z:=(\mathbb N\times\mathbb N)/{\sim}
             =\{[a,b]\mid a,b\in\mathbb N\}.
\]
Bis zum Ende der Konstruktion bezeichnet \(\mathbb Z\) genau dieses
konkrete Modell. Es ist eine Menge, denn jede Klasse ist eine Teilmenge
von \(\mathbb N\times\mathbb N\), und der Quotient ist eine Menge
solcher Teilmengen. Die Konstruktion ist damit durch vorhandene
mengentheoretische Begriffe beschrieben.

\IntegersPairGrid

Die entscheidende Rechenregel für die neuen Objekte lautet
\[
  [a,b]=[c,d]
  \quad\Longleftrightarrow\quad
  a+d=b+c.
\]
Sind die Paare äquivalent, so zeigt die Transitivität, dass jedes
Element der ersten Klasse auch zur zweiten gehört. Die Symmetrie
liefert die umgekehrte Inklusion. Die Klassen sind somit gleich.
Sind umgekehrt die Klassen gleich, so gehört \((a,b)\) wegen der
Reflexivität zur eigenen und deshalb auch zur anderen Klasse. Dies
bedeutet genau \((a,b)\sim(c,d)\), also die Kreuzsummengleichung.
So entsteht \FormulaRefAuto{IntegerClassEquality}[thm].

Insbesondere sind \([2,5]\), \([3,6]\) und \([0,3]\) dasselbe
Objekt. Ebenso gilt für jedes \(n\in\mathbb N\)
\[
  [n,n]=[0,0],
\]
weil \(n+0=n+0\). Wir haben damit Gleichheit und eine Menge neuer
Objekte. Rechnen können wir mit ihnen erst, wenn wir Operationen auf
diesen Klassen definiert haben.

\par\Needspace{10\baselineskip}
\hypertarget{integers.reading.embedding}{}
\section*{3. Wo bleiben die natürlichen Zahlen?}

Die neue Zahl \([n,0]\) soll die bisherige natürliche Zahl \(n\)
wiedergeben. Wir definieren daher die Abbildung
\[
  \iota\colon\mathbb N\longrightarrow\mathbb Z,
  \qquad \iota(n):=[n,0].
\]
Sie ist wohldefiniert, denn jedes Paar \((n,0)\) besitzt eine Klasse
im Quotienten. Außerdem ist sie injektiv. Aus
\(\iota(m)=\iota(n)\) folgt nämlich
\[
  [m,0]=[n,0]
  \quad\Longrightarrow\quad
  m+0=0+n
  \quad\Longrightarrow\quad m=n.
\]
Verschiedene natürliche Zahlen bleiben also verschieden. Dies ist
\FormulaRefAuto{NaturalIntegerEmbeddingInjective}[thm].

Die ursprüngliche natürliche Zahl \(n\) und die Klasse \([n,0]\)
sind dabei zunächst verschiedene Objekte. Wir haben nicht bewiesen,
dass die ursprüngliche Menge \(\mathbb N\) wörtlich eine Teilmenge
des Quotienten ist. Eine Teilmenge von \(\mathbb Z\) ist vielmehr
das Bild
\[
  \iota[\mathbb N]=\{[n,0]\mid n\in\mathbb N\}.
\]
Zwischen \(\mathbb N\) und dieser Teilmengenkopie ist \(\iota\)
eine Bijektion: Die Injektivität wurde eben gezeigt; jedes Element
des Bildes wird nach dessen Definition getroffen. Wenn wir später
\(\mathbb N\subseteq\mathbb Z\) schreiben, ist diese durch
\(\iota\) identifizierte Kopie gemeint.

Für den Anfang halten wir die Unterscheidung sichtbar und setzen
\[
  0_{\mathbb Z}:=\iota(0)=[0,0],
  \qquad
  1_{\mathbb Z}:=\iota(1)=[1,0].
\]
In einem Kontext ganzer Zahlen wird später das Zeichen \(n\) kurz
für \(\iota(n)\) stehen. Einträge innerhalb von \([a,b]\) bleiben
stets natürliche Zahlen. Die Einbettung sorgt dafür, dass die gleiche
Ziffernschreibweise die beiden Zahlbereiche verständlich verbindet.

Noch ist nur die Erhaltung der Gleichheit gesichert. Dass die
Einbettung auch Addition, Multiplikation und Ordnung erhält, können
wir erst nach deren Konstruktion beweisen. Eine Zahlbereichserweiterung
soll schließlich nicht bloß die alten Zahlen aufzählen. Sie soll die
bereits bewiesenen Rechnungen mit ihnen fortsetzen.

\par\Needspace{10\baselineskip}
\hypertarget{integers.reading.addition}{}
\section*{4. Negation, Addition und Subtraktion}

Bei der gedachten Differenz \(a-b\) kehrt das Vertauschen der
Einträge das Vorzeichen um. Für die Klassen liegt daher die Vorschrift
\[
  -[a,b]:=[b,a]
\]
nahe. Ob sie eine Operation auf Zahlen definiert, hängt von einer
entscheidenden Prüfung ab: Verschiedene Repräsentanten derselben
Eingabe müssen dieselbe Ausgabe liefern. Andernfalls würde das
Ergebnis von unserer Schreibweise abhängen.

\medskip\noindent\textbf{Dieselbe Zahl muss dasselbe Gegenstück erhalten.}
Wir nehmen an, dass \([a,b]=[a',b']\) gilt. Zu zeigen ist, dass
auch nach dem Vertauschen dieselbe Klasse entsteht, also
\([b,a]=[b',a']\). Das Gleichheitskriterium aus Abschnitt~2
übersetzt Voraussetzung und Ziel in natürliche Gleichungen:
\[
  \begin{aligned}
    [a,b]=[a',b']&\quad\Longleftrightarrow\quad a+b'=b+a',\\
    [b,a]=[b',a']&\quad\Longleftrightarrow\quad b+a'=a+b'.
  \end{aligned}
\]
Die zweite Gleichung ist die erste mit vertauschten Seiten. Daher
folgt das Ziel aus der Voraussetzung. Damit hängt die Negation
nicht vom gewählten Repräsentanten ab.

Zwei aufeinanderfolgende Negationen vertauschen zweimal und liefern
\[
  -(-[a,b])=[a,b].
\]
Für die Null kann man dies direkt an jedem Paar mit gleichen
Einträgen sehen: Nach Abschnitt~2 ist \(0_{\mathbb Z}=[a,a]\)
für jedes \(a\in\mathbb N\), und Vertauschen ändert dieses Paar
nicht. Also gilt
\[
  -0_{\mathbb Z}=-[a,a]=[a,a]=0_{\mathbb Z}.
\]

\medskip\noindent\textbf{Eine Summe aus zwei Paaren.}
Für die Addition verwenden wir die Vorschrift
\[
  [a,b]+[c,d]:=[a+c,b+d].
\]
Sie addiert die jeweils ersten und die jeweils zweiten Einträge.
Auch hier muss das Ergebnis bei einem Wechsel der Repräsentanten
gleich bleiben. Wir nehmen deshalb an, dass
\[
  [a,b]=[a',b'],\qquad [c,d]=[c',d'].
\]
Unser Ziel ist die Gleichheit der beiden vorgeschlagenen Summen:
\[
  [a+c,b+d]=[a'+c',b'+d'].
\]
Die Voraussetzungen lauten als Kreuzsummengleichungen
\[
  a+b'=b+a',
  \qquad c+d'=d+c'.
\]
Für das Ziel brauchen wir die Kreuzsummengleichung
\((a+c)+(b'+d')=(b+d)+(a'+c')\). Sie folgt durch Umordnen
und Einsetzen der beiden Voraussetzungen:
\[
  \begin{aligned}
    (a+c)+(b'+d')
      &=(a+b')+(c+d')\\
      &=(b+a')+(d+c')\\
      &=(b+d)+(a'+c').
  \end{aligned}
\]
Dies ist die Kreuzsummengleichung für
\([a+c,b+d]=[a'+c',b'+d']\).
Die Vorschrift definiert also eine eindeutige Summe zweier Klassen.
Die formalen Prüfungen heißen
\FormulaRefAuto{IntegerNegationCompatible}[thm] und
\FormulaRefAuto{IntegerAdditionCompatible}[thm].

Die Ergebnisse liegen wieder in \(\mathbb Z\), weil alle neuen
Paare natürliche Einträge haben. Damit ist die Abgeschlossenheit
unter Negation und Addition gesichert. Nun können wir die
Rechengesetze an beliebigen Repräsentanten prüfen. Gerade die eben
bewiesene Unabhängigkeit erlaubt uns diese freie Wahl.

\medskip\noindent\textbf{Die Klammerung einer Summe ist frei.}
Wir wollen das Assoziativgesetz
\((z+w)+u=z+(w+u)\) zeigen. Für \(z=[a,b]\),
\(w=[c,d]\) und \(u=[e,f]\) liefert die Additionsdefinition
\[
  \begin{aligned}
    (z+w)+u&=[(a+c)+e,(b+d)+f],\\
    z+(w+u)&=[a+(c+e),b+(d+f)].
  \end{aligned}
\]
Wegen der natürlichen Assoziativität ist
\((a+c)+e=a+(c+e)\) und ebenso \((b+d)+f=b+(d+f)\).
Die beiden Paare haben also dieselben Einträge und damit dieselbe
Klasse. Genau dies ist die gewünschte Gleichheit der beiden
geklammerten Summen.

\medskip\noindent\textbf{Die Reihenfolge einer Summe ist frei.}
Das Kommutativgesetz \(z+w=w+z\) folgt ebenso aus der
natürlichen Kommutativität:
\[
  z+w=[a+c,b+d]=[c+a,d+b]=w+z.
\]

\medskip\noindent\textbf{Null ist neutral; das Gegenstück hebt eine Zahl auf.}
Ein neutrales Element der Addition lässt jede Zahl beim Addieren
unverändert. Für \(0_{\mathbb Z}=[0,0]\) erhalten wir
\[
  [a,b]+[0,0]=[a+0,b+0]=[a,b].
\]
Das Gegenstück \(-[a,b]=[b,a]\) soll dagegen mit der ursprünglichen
Zahl zusammen null ergeben. Tatsächlich ist
\[
  [a,b]+[b,a]=[a+b,b+a]=[a+b,a+b]=[0,0].
\]
Im vorletzten Schritt benutzen wir \(b+a=a+b\), im letzten
die bereits bekannte Regel \([n,n]=[0,0]\).
Damit besitzt jede ganze Zahl ein additives Gegenstück. Die gleichen
Gesetze mit vertauschter Reihenfolge folgen aus der Kommutativität.

\medskip\noindent\textbf{Jetzt darf subtrahiert werden.}
Wir definieren
\[
  z-w:=z+(-w).
\]
Subtraktion ist also eine abgekürzte Addition mit einem Gegenstück.
In Paaren lautet die Formel
\[
  \begin{aligned}
    [a,b]-[c,d]
      &=[a,b]+\bigl(-[c,d]\bigr)\\
      &=[a,b]+[d,c]\\
      &=[a+d,b+c].
  \end{aligned}
\]
Die drei Schritte sind die Definition der Subtraktion, die
Definition der Negation und die Definition der Addition.
Damit ist auch unser Ausgangsproblem lösbar. Es gilt
\[
  [0,3]+[5,0]=[5,3]=[2,0],
\]
denn \(5+0=3+2\). In der eingeführten Ziffernschreibweise heißt
dies \(-3+5=2\). Die Lösung \(-3\) ist nun die Klasse \([0,3]\)
und nicht mehr bloß ein zusätzliches Zeichen.

Auch die frühere Addition bleibt erhalten:
\[
  \iota(m)+\iota(n)
    =[m,0]+[n,0]
    =[m+n,0]
    =\iota(m+n).
\]
Dies ist \FormulaRefAuto{NaturalIntegerEmbeddingAdditive}[thm].
Jede ganze Zahl lässt sich außerdem als Differenz eingebetteter
natürlicher Zahlen lesen:
\[
  [a,b]=\iota(a)-\iota(b).
\]
Diese Gleichung ist jetzt eine bewiesene Aussage über definierte
Operationen. Die anfängliche Deutung eines Paares als Differenz ist
damit mathematisch eingelöst.

Die bisher bewiesenen Gesetze machen \((\mathbb Z,+,0_{\mathbb Z})\)
zu einer abelschen Gruppe. Für die weitere Entwicklung müssen wir
diesen Namen nicht voraussetzen: Gemeint sind genau die
Assoziativität, Kommutativität, Nullgesetze und Gegenstücke, die wir
soeben hergeleitet haben. Die Gruppen-Lesefassung in Band~41 stellt
später diese additive Seite in einen allgemeineren Zusammenhang.

\par\Needspace{10\baselineskip}
\hypertarget{integers.reading.multiplication}{}
\section*{5. Wie die Multiplikation entsteht}

Die Addition allein genügt noch nicht für die Ganzzahlarithmetik.
Für die Multiplikation soll die Vorstellung einer Differenz ebenfalls
passen. Die vertraute Rechnung
\[
  (a-b)(c-d)=(ac+bd)-(ad+bc)
\]
legt eine Vorschrift nahe. Diese Rechnung ist zunächst nur eine
Motivation: Wir verwenden die noch zu begründenden Ganzzahlgesetze
nicht als Beweis. Die tatsächliche Definition soll durch die
natürlichen Einträge des Paares
\[
  [a,b]\cdot[c,d]:=[ac+bd,ad+bc]
\]
gegeben werden. Vor ihrer Verwendung prüfen wir die
Repräsentantenunabhängigkeit.

\medskip\noindent\textbf{Wechsel des ersten Repräsentanten.}
Wir halten das zweite Paar \((c,d)\) fest und nehmen
\([a,b]=[a',b']\) an. Obwohl wir für den ersten Faktor ein
anderes Paar einsetzen, soll dieselbe Produktklasse entstehen.
Zu zeigen ist also
\[
  [ac+bd,ad+bc]=[a'c+b'd,a'd+b'c].
\]
Nach dem Klassengleichheitskriterium genügt dafür die natürliche
Gleichung
\[
  (ac+bd)+(a'd+b'c)=(ad+bc)+(a'c+b'd).
\]
Aus der Voraussetzung folgt \(a+b'=b+a'\). Multiplikation
dieser Gleichung mit \(c\) beziehungsweise mit \(d\)
und die natürliche Distributivität liefern
\[
  ac+b'c=bc+a'c,
  \qquad ad+b'd=bd+a'd.
\]
Wir lesen die zweite Gleichung rückwärts und setzen beide ein:
\[
  \begin{aligned}
    (ac+bd)+(a'd+b'c)
      &=(ac+b'c)+(bd+a'd)\\
      &=(bc+a'c)+(ad+b'd)\\
      &=(ad+bc)+(a'c+b'd).
  \end{aligned}
\]
Dies ist genau die benötigte Kreuzsummengleichung. Die beiden
Produktklassen sind also gleich. Bei festem zweiten Paar ist der
Wechsel des ersten damit erlaubt.

Für den Wechsel des zweiten Paares benutzen wir die Symmetrie der
Paarvorschrift. Schon in \(\mathbb N\) gilt
\[
  (ac+bd,ad+bc)=(ca+db,cb+da).
\]
Wir können daher die Rollen der beiden Paare vertauschen, das gerade
gezeigte Argument auf das nun erste Paar anwenden und zurücktauschen.
Ein Wechsel beider Repräsentanten wird in diese beiden einzelnen
Wechsel zerlegt. Die beiden resultierenden Klassengleichheiten werden
durch Transitivität verbunden. So ist die Produktvorschrift unabhängig
von beiden Darstellungen; dies ist
\FormulaRefAuto{IntegerMultiplicationCompatible}[thm].

Alle Einträge des Produktpaares sind natürliche Zahlen. Daher ist
auch die Multiplikation auf ganz \(\mathbb Z\) definiert und
liefert wieder eine ganze Zahl. Erst ab jetzt dürfen wir die neue
Multiplikation ohne Hinweis auf Repräsentanten verwenden.

\medskip\noindent\textbf{Die Rechengesetze des Produkts.}
Die eben verwendete Symmetrie der Paarvorschrift beweist zugleich
\(z\cdot w=w\cdot z\). Für die Eins ergibt sich
\[
  [a,b]\cdot[1,0]=[a\cdot1+b\cdot0,a\cdot0+b\cdot1]=[a,b].
\]
Wegen der Kommutativität ist \([1,0]\) auch auf der anderen Seite
neutral. Ferner liefert das Einsetzen von \([0,0]\) das Produkt
\([0,0]\); somit gilt \(z\cdot0_{\mathbb Z}=0_{\mathbb Z}\).

\medskip\noindent\textbf{Assoziativität der Multiplikation.}
Wir wollen zeigen, dass auch bei drei Faktoren die Klammerung
keinen Einfluss hat:
\[
  (z\cdot w)\cdot u=z\cdot(w\cdot u).
\]
Dazu setzen wir wieder
\(z=[a,b]\), \(w=[c,d]\) und \(u=[e,f]\).
Beim Produkt \((z\cdot w)\cdot u\) sind die Einträge nach
natürlichem Ausmultiplizieren
\[
  \begin{aligned}
    (ac+bd)e+(ad+bc)f&=ace+bde+adf+bcf,\\
    (ac+bd)f+(ad+bc)e&=acf+bdf+ade+bce.
  \end{aligned}
\]
Für die andere Klammerung \(z\cdot(w\cdot u)\) erhalten wir
\[
  \begin{aligned}
    a(ce+df)+b(cf+de)&=ace+adf+bcf+bde,\\
    a(cf+de)+b(ce+df)&=acf+ade+bce+bdf.
  \end{aligned}
\]
Die ersten Einträge unterscheiden sich nur in der Reihenfolge
ihrer Summanden, ebenso die zweiten. Die Paare sind gleich,
also auch ihre Klassen. Das beweist
\((z\cdot w)\cdot u=z\cdot(w\cdot u)\).

\medskip\noindent\textbf{Distributivität: Eine Summe wird gliedweise multipliziert.}
Unser nächstes Ziel ist
\[
  z\cdot(w+u)=z\cdot w+z\cdot u.
\]
Wir verwenden weiterhin \(z=[a,b]\), \(w=[c,d]\) und
\(u=[e,f]\). Links bilden wir zuerst die Summe
\(w+u=[c+e,d+f]\) und danach das Produkt mit \(z\):
\[
  z\cdot(w+u)
    =\bigl[a(c+e)+b(d+f),\ a(d+f)+b(c+e)\bigr].
\]
Rechts bilden wir zuerst die beiden Produkte
\[
  z\cdot w=[ac+bd,ad+bc],\qquad
  z\cdot u=[ae+bf,af+be]
\]
und addieren anschließend:
\[
  z\cdot w+z\cdot u
    =\bigl[(ac+bd)+(ae+bf),\ (ad+bc)+(af+be)\bigr].
\]
Nun vergleichen wir die Einträge. Die natürliche Distributivität
erlaubt das Ausmultiplizieren; natürliche Assoziativität und
Kommutativität erlauben das Umordnen:
\[
  \begin{aligned}
    a(c+e)+b(d+f)
      &=ac+ae+bd+bf\\
      &=(ac+bd)+(ae+bf),\\[.3em]
    a(d+f)+b(c+e)
      &=ad+af+bc+be\\
      &=(ad+bc)+(af+be).
  \end{aligned}
\]
Somit sind die ersten Einträge gleich und ebenso die zweiten.
Beide Seiten unserer Zielgleichung sind dieselbe Klasse. Das
beweist die Linksdistributivität. Die Rechtsdistributivität
\((z+w)\cdot u=z\cdot u+w\cdot u\) folgt anschließend aus der
Kommutativität:
\[
  (z+w)\cdot u
    =u\cdot(z+w)
    =u\cdot z+u\cdot w
    =z\cdot u+w\cdot u.
\]
Die ausführlichen Ableitungen stehen bei
\FormulaRefAuto{IntMulAssociative}[thm] und
\FormulaRefAuto{IntMulLeftDistributive}[thm]. Wir haben damit alle
elementaren Rechengesetze der Ganzzahlarithmetik am Modell geprüft.
Zusammen beschreiben sie einen kommutativen Ring mit Eins.

\medskip\noindent\textbf{Vorzeichenregeln werden hergeleitet.}
Warum wird aus dem Produkt zweier negativer Zahlen ein positives?
Die Antwort ist bereits in der Produktdefinition enthalten. Negation
vertauscht die Einträge. Deshalb gilt für beliebige \(z=[a,b]\)
und \(w=[c,d]\)
\[
  \begin{aligned}
    (-z)\cdot w
      &=[bc+ad,bd+ac]\\
      &=[ad+bc,ac+bd]\\
      &=-(z\cdot w).
  \end{aligned}
\]
Mit der Kommutativität ergibt sich auch
\(z\cdot(-w)=-(z\cdot w)\). Zweimaliges Anwenden und die
doppelte Negation liefern
\[
  (-z)\cdot(-w)
    =-\bigl(z\cdot(-w)\bigr)
    =-\bigl(-(z\cdot w)\bigr)
    =z\cdot w.
\]
Die Vorzeichenregel ist also keine zusätzliche freie Festsetzung.
Sie folgt aus der konstruierten Multiplikation. Band~17 zeigt dieselben
Regeln auch aus den bereits hergeleiteten arithmetischen Axiomen, etwa
in \FormulaRefAuto{IntegerNegativeTimesNegative}[thm].

Die natürliche Multiplikation bleibt ebenfalls erhalten:
\[
  \iota(m)\cdot\iota(n)
    =[m,0]\cdot[n,0]
    =[mn,0]
    =\iota(mn).
\]
Damit ist \FormulaRefAuto{NaturalIntegerEmbeddingMultiplicative}[thm]
begründet. Beispielsweise folgt nun
\[
  (-\iota(2))\cdot(-\iota(3))=\iota(6),
  \qquad
  \iota(2)\cdot(-\iota(3))=-\iota(6).
\]
In der üblichen Schreibweise sind das \((-2)\cdot(-3)=6\) und
\(2\cdot(-3)=-6\). Die eingebetteten natürlichen Zahlen rechnen
also wie zuvor, während ihre Gegenstücke den neuen negativen Teil
des Zahlbereichs erschließen.

\par\Needspace{10\baselineskip}
\hypertarget{integers.reading.normalforms}{}
\section*{6. Jede ganze Zahl besitzt eine Vorzeichenform}

Ein Paar kann sehr viele Darstellungen derselben Zahl liefern. Für
das Verständnis der neuen Menge ist deshalb eine einfachere Form
wichtig: Jede Klasse lässt sich als eingebettete natürliche Zahl
oder als deren Gegenstück schreiben.

Sei \(z=[a,b]\). Die natürlichen Zahlen \(a\) und \(b\) sind
vergleichbar. Gilt \(b\leq a\), so gibt es nach der natürlichen
Abstandsbeschreibung ein \(n\in\mathbb N\) mit \(b+n=a\).
Dann ist
\[
  a+0=b+n,
  \qquad [a,b]=[n,0]=\iota(n).
\]
Gilt dagegen \(a\leq b\), wählen wir \(n\in\mathbb N\) mit
\(a+n=b\). Nun ergibt die Kreuzsummengleichung
\[
  a+n=b+0,
  \qquad [a,b]=[0,n]=-\iota(n).
\]
Damit ist in allen Fällen
\[
  z=\iota(n)\quad\text{oder}\quad z=-\iota(n)
  \qquad\text{für ein }n\in\mathbb N.
\]
Dies ist \FormulaRefAuto{IntegerSignedNormalForm}[thm]. Der Beweis
benutzt die natürliche Ordnung; eine Ganzzahlordnung haben wir
noch nicht vorausgesetzt.

Zum Beispiel ergibt \(5+3=8\) die Gleichheit
\([8,5]=[3,0]\), während \(2+5=7\) zu
\([2,7]=[0,5]\) führt. Die vielen Paare erzeugen somit keine
zusätzlichen Arten von Zahlen zwischen den bekannten positiven und
negativen Vorzeichenformen.

Bei Null fallen die beiden Formen zusammen:
\(\iota(0)=-\iota(0)=0_{\mathbb Z}\).
Die Aussage der Normalform verlangt daher keine unterschiedslos
eindeutige Wahl eines Vorzeichens. Sie sichert zunächst die
Existenz einer der beiden Darstellungen. Genau diese Existenz wird
später die zweiseitige Induktion ermöglichen: Der positive und der
negative Strahl zusammen erfassen jede Klasse.

\par\Needspace{10\baselineskip}
\hypertarget{integers.reading.order}{}
\section*{7. Eine Ordnung, die zum Rechnen passt}

Die Konstruktion von Addition und Multiplikation sagt noch nicht,
welche ganze Zahl kleiner ist. Aus der Vorstellung \(a-b\leq c-d\)
liegt erneut ein Vergleich von Kreuzsummen nahe. Wir setzen
\[
  [a,b]\leq[c,d]
  \quad\Longleftrightarrow\quad
  a+d\leq c+b.
\]
Rechts steht die bereits bekannte natürliche Ordnung, links die
neue Ganzzahlordnung. Zuerst müssen wir auch diese Vorschrift
unabhängig von den Repräsentanten machen.

Seien also \([a,b]=[a',b']\) und \([c,d]=[c',d']\).
Dann gelten die natürlichen Gleichungen
\[
  a+b'=b+a',
  \qquad c+d'=d+c'.
\]
Sie geben nach Umordnen
\[
  \begin{aligned}
    (a+d)+(b'+d')&=(a'+d')+(b+d),\\
    (c+b)+(b'+d')&=(c'+b')+(b+d).
  \end{aligned}
\]
Wir dürfen zu beiden Seiten einer natürlichen Ungleichung dieselbe
Zahl addieren und diesen Summanden wieder kürzen. Daher ist
\(a+d\leq c+b\) gleichbedeutend mit der Ungleichung zwischen
den beiden linken Seiten und damit zwischen den beiden rechten
Seiten. Nach Kürzen von \(b+d\) bleibt
\(a'+d'\leq c'+b'\). Jeder Schritt ist eine Äquivalenz;
die Aussage gilt in beiden Richtungen. Dies beweist
\FormulaRefAuto{IntegerOrderCompatible}[thm].

\medskip\noindent\textbf{Warum dies eine totale Ordnung ist.}
Reflexivität folgt aus \(a+b\leq a+b\). Sind
\([a,b]\leq[c,d]\) und \([c,d]\leq[a,b]\), so liefern die
beiden Kreuzsummenvergleiche \(a+d=c+b\), also
\([a,b]=[c,d]\). Das ist die Antisymmetrie.

Für die Transitivität seien
\[
  a+d\leq c+b,
  \qquad c+f\leq e+d.
\]
Durch Addieren natürlicher Zahlen und Umordnen folgt
\[
  (a+f)+d
    \leq(c+b)+f
    =(c+f)+b
    \leq(e+d)+b
    =(e+b)+d.
\]
Das natürliche Ordnungskürzungsgesetz ergibt
\(a+f\leq e+b\), also \([a,b]\leq[e,f]\).
Schließlich sind \(a+d\) und \(c+b\) als natürliche Zahlen
vergleichbar. Somit ist auch jedes Paar ganzer Zahlen vergleichbar.
Zusammen ist dies \FormulaRefAuto{IntegerOrderTotal}[thm].

Die strikte Ordnung definieren wir durch
\[
  z<w\quad\Longleftrightarrow\quad z\leq w\text{ und }z\ne w.
\]
Für die eingebetteten natürlichen Zahlen ergibt sich unmittelbar
\[
  \iota(m)\leq\iota(n)
  \quad\Longleftrightarrow\quad m\leq n.
\]
Wegen der Injektivität gilt dasselbe für den strikten Vergleich.
Insbesondere folgen aus \(0<1\) in \(\mathbb N\) die Aussagen
\(0_{\mathbb Z}<1_{\mathbb Z}\) und
\(0_{\mathbb Z}\ne1_{\mathbb Z}\).

\medskip\noindent\textbf{Addition verschiebt die Ordnung.}
Für \(z=[a,b]\), \(w=[c,d]\) und \(u=[e,f]\) gilt
\[
  \begin{aligned}
    z+u\leq w+u
    &\quad\Longleftrightarrow\quad
       (a+e)+(d+f)\leq(c+e)+(b+f)\\
    &\quad\Longleftrightarrow\quad
       (a+d)+(e+f)\leq(c+b)+(e+f)\\
    &\quad\Longleftrightarrow\quad a+d\leq c+b.
  \end{aligned}
\]
Damit ist
\[
  z\leq w\quad\Longleftrightarrow\quad z+u\leq w+u.
\]
Die Ordnung wird durch Addition derselben ganzen Zahl erhalten und
reflektiert. Das letzte Wort bedeutet, dass man aus der verschobenen
Ungleichung zur ursprünglichen zurückkehren darf. Die formale
Aussage ist \FormulaRefAuto{IntegerOrderTranslationModel}[thm].

Auch Gleichheit wird durch eine gemeinsame Addition erhalten und
reflektiert: Für die Rückrichtung addieren wir das Gegenstück des
gemeinsamen Summanden. Deshalb gilt die Translationsäquivalenz
ebenso für \(<\). Addieren wir zu \(z\leq w\) die Zahl
\(-z-w\), so erhalten wir \(-w\leq-z\); dieselbe Überlegung
rückwärts beweist die Umkehrung. Negation kehrt die Ordnung also um.

\medskip\noindent\textbf{Positive Zahlen und positive Produkte.}
Eine ganze Zahl heißt positiv, wenn \(0_{\mathbb Z}<z\).
Für \(z=[a,b]\) bedeutet das nach der Ordnungsdefinition
\(b\leq a\) und nach dem Gleichheitskriterium \(a\ne b\).
Also ist \(b<a\). Es gibt folglich eine natürliche Zahl \(n>0\)
mit \(b+n=a\), und die Normalform liefert \(z=\iota(n)\).
Umgekehrt ist jedes \(\iota(n)\) mit \(n>0\) positiv, weil
die Einbettung den natürlichen strikten Vergleich erhält.

Sind \(z=\iota(m)\) und \(w=\iota(n)\) positiv, so sind
\(m,n>0\). Das natürliche Produkt \(mn\) ist positiv, und
\[
  z\cdot w=\iota(mn)>0_{\mathbb Z}.
\]
So entsteht die Positivität eines Produkts aus der natürlichen
Multiplikation. Sie wird in
\FormulaRefAuto{IntegerPositiveProductModel}[thm] festgehalten.
Für nichtnegative Faktoren folgt daraus die nichtnegative Variante:
Ist einer null, ist das Produkt null; sind beide ungleich null,
sind beide positiv und auch ihr Produkt ist positiv.

\medskip\noindent\textbf{Warum positive Faktoren gekürzt werden dürfen.}
Seien \(a,c,d\in\mathbb Z\) mit \(d>0_{\mathbb Z}\). Wir zeigen
\[
  a\leq c\quad\Longleftrightarrow\quad a\cdot d\leq c\cdot d.
\]
Durch Translation ist \(a\leq c\) gleichbedeutend mit
\(0_{\mathbb Z}\leq c-a\). Das Produkt \((c-a)\cdot d\)
ist dann nichtnegativ. Distributivität und die Vorzeichenregel ergeben
\[
  (c-a)\cdot d=c\cdot d-a\cdot d.
\]
Nach erneuter Translation folgt \(a\cdot d\leq c\cdot d\).

Für die Rückrichtung nehmen wir \(a\cdot d\leq c\cdot d\)
an. Wäre \(c<a\), so wäre \(a-c>0_{\mathbb Z}\).
Das Produkt \((a-c)\cdot d\) wäre strikt positiv, also
\(c\cdot d<a\cdot d\). Das widerspricht der angenommenen
Ungleichung. Da die Ordnung total ist, bleibt \(a\leq c\).
Damit ist die Faktorordnung aus
\FormulaRefAuto{IntegerPositiveFactorOrderModel}[thm] bewiesen.

Aus ihr folgt auch das Kürzen eines positiven Faktors in einer
Gleichung: Aus \(a\cdot d=c\cdot d\) erhalten wir beide
Ungleichungen \(a\leq c\) und \(c\leq a\), also \(a=c\).
Zusammen mit der Erhaltung von Ungleichungen folgt die Erhaltung
des strikten Vergleichs. Die Voraussetzung \(d>0_{\mathbb Z}\)
ist dabei wesentlich. Der Faktor null unterscheidet keine Zahlen;
bei negativen Faktoren kehrt sich die Richtung um.

Die Ordnung ist damit vollständig mit den Rechenoperationen
verbunden: Addition verschiebt sie, Negation kehrt sie um, und
Multiplikation mit einer positiven Zahl erhält und reflektiert sie.
Diese Eigenschaften bilden mit der Ringarithmetik den geordneten
Teil der späteren Axiomatik.

\par\Needspace{10\baselineskip}
\hypertarget{integers.reading.induction}{}
\section*{8. Schritte um eins und zweiseitige Induktion}

In \(\mathbb Z\) können wir von jeder Zahl aus sowohl eins
addieren als auch eins subtrahieren. Auf Repräsentanten lauten die
beiden Schritte
\[
  [a,b]+1_{\mathbb Z}=[a+1,b],
  \qquad
  [a,b]-1_{\mathbb Z}=[a,b+1].
\]
Für die zweite Formel benutzen wir
\(-1_{\mathbb Z}=[0,1]\). Beide Schritte sind überall definiert,
weil die Addition und die Subtraktion abgeschlossen sind.

Sie machen einander rückgängig. Aus den Additionsgesetzen folgt
\[
  \begin{aligned}
    (z+1_{\mathbb Z})-1_{\mathbb Z}
       &=z+(1_{\mathbb Z}+(-1_{\mathbb Z}))=z,\\
    (z-1_{\mathbb Z})+1_{\mathbb Z}
       &=z+((-1_{\mathbb Z})+1_{\mathbb Z})=z.
  \end{aligned}
\]
Die Abbildungen \(z\mapsto z+1_{\mathbb Z}\) und
\(z\mapsto z-1_{\mathbb Z}\) sind daher gegenseitige
Umkehrabbildungen. Anders als der natürliche Nachfolger trifft
der Ganzzahlnachfolger auch die Null: Ihr Vorgänger ist
\(-1_{\mathbb Z}\).

Die positive Orientierung eines Schrittes folgt aus der
Translationsverträglichkeit der strikten Ordnung. Da
\(0_{\mathbb Z}<1_{\mathbb Z}\), gilt für jedes \(z\)
\[
  z<z+1_{\mathbb Z}.
\]
Dies ist \FormulaRefAuto{IntegerSuccessorPositive}[thm]. Wiederholte
Vorwärtsschritte können deshalb keine Schleife bilden. Auf den
Vorzeichenformen haben die beiden Richtungen eine besonders einfache
Beschreibung:
\[
  \begin{aligned}
    \iota(n)+1_{\mathbb Z}&=\iota(n+1),\\
    -\iota(n)-1_{\mathbb Z}&=-\iota(n+1).
  \end{aligned}
\]
Die erste Gleichung ist die Additivität der Einbettung; die zweite
folgt unmittelbar aus \([0,n]+[0,1]=[0,n+1]\).

\IntegersTwoRays

\medskip\noindent\textbf{Positive Schritte liefern Schranken.}
Die Normalform erklärt auch eine im Skript später benötigte
Abschätzung. Ist \(z=\iota(m)\), so gilt
\(z<\iota(m+1)\) nach der positiven Orientierung. Ist
\(z=-\iota(m)\), so gilt
\(z\leq0_{\mathbb Z}<\iota(m+1)\). In beiden Fällen finden
wir also eine eingebettete natürliche Zahl oberhalb von \(z\).

Eine positive ganze Zahl \(b\) hat die Form \(\iota(n)\)
mit natürlichem \(n>0\). Aus \(1\leq n\) folgt
\(1_{\mathbb Z}\leq b\). Da \(\iota(m+1)>0_{\mathbb Z}\),
können wir diese Ungleichung mit \(\iota(m+1)\) multiplizieren.
Zusammen erhalten wir
\[
  z<\iota(m+1)\leq\iota(m+1)\cdot b.
\]
Damit können auch natürliche Vielfache jedes positiven Faktors
als obere Schranken verwendet werden. Dies ist der Gedanke hinter
\FormulaRefAuto{RationalIntegerFractionBoundModel}[thm]. Er zeigt,
wie Normalform, Schrittstruktur und Faktorordnung zusammenwirken.

\medskip\noindent\textbf{Zwei natürliche Induktionen reichen aus.}
Sei \(P(z)\) eine Aussage über ganze Zahlen. Wir nehmen an:
\[
  \begin{gathered}
    P(0_{\mathbb Z}),\\
    \forall z\in\mathbb Z\,
      \bigl(P(z)\Longrightarrow P(z+1_{\mathbb Z})\bigr),\\
    \forall z\in\mathbb Z\,
      \bigl(P(z)\Longrightarrow P(z-1_{\mathbb Z})\bigr).
  \end{gathered}
\]
Wir wollen zeigen, dass \(P\) für jede ganze Zahl gilt.
Der Beweis benutzt zweimal die bereits bekannte Induktion auf
\(\mathbb N\).

Zuerst betrachten wir die Aussage \(P(\iota(n))\). Für \(n=0\)
ist sie die Anfangsvoraussetzung. Gilt \(P(\iota(n))\), so liefert
der Vorwärtsschritt \(P(\iota(n)+1_{\mathbb Z})\), und die
eben bewiesene Gleichung macht daraus \(P(\iota(n+1))\).
Natürliche Induktion ergibt somit
\[
  \forall n\in\mathbb N\,P(\iota(n)).
\]

Für den negativen Strahl betrachten wir \(P(-\iota(n))\).
Bei \(n=0\) benutzen wir
\(-\iota(0)=0_{\mathbb Z}\), also wieder die
Anfangsvoraussetzung. Aus \(P(-\iota(n))\) folgt mit dem
Rückwärtsschritt \(P(-\iota(n)-1_{\mathbb Z})\) und damit
\(P(-\iota(n+1))\). Die zweite natürliche Induktion liefert
\[
  \forall n\in\mathbb N\,P(-\iota(n)).
\]

Nun sei \(z\in\mathbb Z\) beliebig. Nach der Normalform ist
\(z=\iota(n)\) oder \(z=-\iota(n)\) für ein natürliches
\(n\). Im ersten Fall gilt \(P(z)\) durch den ersten, im zweiten
durch den zweiten Induktionslauf. Somit gilt
\[
  \forall z\in\mathbb Z\,P(z).
\]
Dies ist der Modellbeweis von
\FormulaRefAuto{IntegerTwoSidedInductionModel}[thm]. Die neue
Induktionsregel wurde also aus der natürlichen Induktion und der
Normalform gewonnen. Ein Schritt von null nur in die positive
Richtung würde den negativen Strahl nicht erfassen. Genau deshalb
braucht die Ganzzahlinduktion beide Schrittrichtungen.

\par\Needspace{10\baselineskip}
\hypertarget{integers.reading.axioms}{}
\section*{9. Vom konstruierten Modell zur Axiomatik}

Bis hierher war jede ganze Zahl eine Klasse von Paaren. Wir haben
Negation, Addition, Multiplikation und Ordnung auf diesen Klassen
definiert und jede Definition auf Repräsentantenunabhängigkeit
geprüft. Anschließend wurden ihre Gesetze am Modell bewiesen.
Wir können nun zusammenstellen, was für die weitere Verwendung
der ganzen Zahlen tatsächlich benötigt wird.

\medskip\noindent\textbf{Träger, Konstanten und Rechenoperationen.}
Es gibt eine Menge \(\mathbb Z\) mit den beiden ausgezeichneten Elementen \(0\) und
\(1\), einer Negation und zwei zweistelligen Operationen \(+\)
und \(\cdot\). Hier lassen wir die bisher zur Unterscheidung
verwendeten Trägerindizes an \(0_{\mathbb Z}\) und
\(1_{\mathbb Z}\) weg: Alle folgenden Variablen und Konstanten
gehören zum ganzzahligen Kontext. Die Operationen bleiben in
\(\mathbb Z\), und es gelten
\[
  \begin{gathered}
    (z+w)+u=z+(w+u),\\
    z+w=w+z,\\
    z+0=0+z=z,\\
    z+(-z)=(-z)+z=0,
  \end{gathered}
\]
sowie
\[
  \begin{gathered}
    (z\cdot w)\cdot u=z\cdot(w\cdot u),\\
    z\cdot w=w\cdot z,\\
    z\cdot1=1\cdot z=z,\\
    z\cdot(w+u)=z\cdot w+z\cdot u,\\
    (z+w)\cdot u=z\cdot u+w\cdot u.
  \end{gathered}
\]
Diese Regeln sind in
\FormulaRefAuto{IntegerArithmeticAxioms}[def] gesammelt.
Die Bezeichnung kommutativer Ring mit Eins fasst genau diesen
arithmetischen Teil zusammen.

\medskip\noindent\textbf{Ordnung und ihre Verträglichkeit.}
Die Relation \(\leq\) ist reflexiv, antisymmetrisch, transitiv
und total. Dazu kommen
\[
  \begin{gathered}
    z\leq w\quad\Longleftrightarrow\quad z+u\leq w+u,\\
    0\ne1,\qquad 0<1,\\
    0<z\text{ und }0<w\quad\Longrightarrow\quad0<z\cdot w.
  \end{gathered}
\]
Die daraus hergeleitete positive Faktorordnung steht ebenfalls
als benannte Regel zur Verfügung:
\[
  0<d\quad\Longrightarrow\quad
  \bigl(a\leq c\ \Longleftrightarrow\ a\cdot d\leq c\cdot d\bigr).
\]
Diese Zusammenstellung heißt
\FormulaRefAuto{IntegerOrderedArithmeticAxioms}[def]. Die getrennt
aufgeführten Ordnungsregeln entfalten den Begriff der totalen
Ordnung; sie sind keine weiteren unabhängigen Anforderungen.
Ebenso ist die positive Faktorordnung bereits eine Folgerung
der übrigen Gesetze. Ihre eigene Benennung erleichtert spätere
Verweise.

\medskip\noindent\textbf{Schritte und Induktion.}
Die Ausdrücke \(z+1\) und \(z-1\) definieren Abbildungen von
\(\mathbb Z\) nach \(\mathbb Z\). Dabei bleibt
\(z-w=z+(-w)\) die Bedeutung der Subtraktion. Es gelten
\[
  (z+1)-1=z,\qquad (z-1)+1=z,\qquad z<z+1.
\]
Die Schrittabbildungen sind somit keine zusätzlichen
Grundoperationen, sondern durch die Arithmetik bestimmt. Ergänzt
werden sie durch die zweiseitige Induktionsregel: Eine Aussage,
die bei null gilt und bei beiden Schritten erhalten bleibt,
gilt auf ganz \(\mathbb Z\).

Erst die Verbindung dieser drei Teile ist die vollständige
Axiomatik aus \FormulaRefAuto{IntegerPeanoAxioms}[def]. Ein
geordneter Ring allein würde die Rolle der Schritte und der
Induktion noch nicht ausdrücken. Umgekehrt würden zwei
Umkehrabbildungen allein noch keine Multiplikation mit ihren
Gesetzen liefern. Der Modellnachweis muss alle Teile erfassen.
Unsere Konstruktion tut dies; die formale Schlussaussage ist
\FormulaRefAuto{IntQuotientModelsIntegerPeano}[thm].

\medskip\noindent\textbf{Was sich am Abstraktionsschnitt ändert.}
Vor diesem Punkt durfte ein Argument mit \([a,b]\) beginnen
und durch eine Rechnung in \(\mathbb N\) geführt werden.
Damit haben wir nachgewiesen, dass die verlangte Struktur
aus den vorausgesetzten natürlichen Zahlen konstruiert werden
kann. Nach dem Abstraktionsschnitt bildet die vollständige
Axiomatik die Grundlage weiterer Beweise. Die konkrete Klasse
eines Paares ist dann kein Bestandteil der verwendeten Sprache
mehr.

Das hat eine praktische Folge: Wenn wir später etwa ein
Distributivgesetz oder eine Ordnungsverschiebung benötigen,
greifen wir auf die entsprechende Regel zurück. Wir wählen
dafür keine Paarrepräsentanten. Auch die Vorzeichenregeln
können rein aus den arithmetischen Gesetzen hergeleitet werden.
Die Konstruktion erklärt, weshalb das ganze Regelwerk ein
Modell besitzt; die weitere Arithmetik arbeitet mit diesem
Regelwerk selbst.

Bei der Induktion ist dabei eine logische Unterscheidung zu
beachten. Als erststufiges Schema wird sie für die in der
Sprache formulierbaren Aussagen verwendet. Daraus folgt keine
Behauptung, dass alle Modelle des Schemas bis auf Isomorphie
gleich wären. Eine volle zweitstufige Lesart lässt dagegen
beliebige Teilmengen beziehungsweise Eigenschaften des Trägers
zu und ist stärker. Für den hier geführten Modellnachweis
genügen die jeweils vorausgesetzten natürlichen
Induktionsinstanzen. Wir leiten aus ihm keine zusätzliche
Eindeutigkeitsaussage über Modelle ab.

\medskip\noindent\textbf{Der Anschluss an die formalen Differenzen.}
Der Weg von Paaren über eine Kreuzsummenrelation zur Addition
von Klassen ist zugleich der Ausgangspunkt der
\DifferencesReadingLink{} in Band~41. Dort wird die
natürliche Addition durch eine beliebige kommutative kürzbare
Monoidoperation ersetzt. Untersucht wird, weshalb dieselbe
Konstruktion eine abelsche Gruppe und eine injektive Einbettung
liefert.

Die Multiplikation, die Ordnung und die Induktion dieser
Lesefassung beruhen dagegen auf weiteren Eigenschaften der
natürlichen Zahlen. Sie werden nicht allein aus einer
Monoidoperation gewonnen. Damit haben die beiden Texte
verschiedene Ziele: Band~17 entwickelt die ganzen Zahlen bis
zur vollständigen Axiomatik; Band~41 erklärt die allgemeine
additive Methode hinter dem ersten Teil dieser Konstruktion.

\clearpage
\hypertarget{integers.reading.history}{}
\section*{10. Zur Geschichte der ganzen Zahlen}

Die Rechenregeln für negative Zahlen haben eine lange Geschichte.
Ein wichtiger früher Beleg ist Brahmaguptas
\emph{Brahmasphutasiddhanta} aus dem Jahr~628. Dort begegnen
positive und negative Größen in der Sprache von Vermögen und
Schulden. Brahmagupta beschreibt, wie man solche Größen addiert
und subtrahiert, und nennt bereits die Regel, dass das Produkt
zweier negativer Größen positiv ist.\footnote{Brahmagupta,
\emph{Brahmasphutasiddhanta}, Rechenregeln in der englischen
Übersetzung von H.~T.~Colebrooke (1817),
\href{https://www.hist-math.fr/textes/Colebrooke1817_BrahmeguptaBhaskaracharya.pdf}{\emph{Algebra, with Arithmetic and Mensuration, from the Sanscrit of Brahmegupta and Bháscara}},
S.~339, Abschnitt~II, \emph{Algorithm}.
Zur Datierung: J.~J.~O'Connor und E.~F.~Robertson,
\href{https://mathshistory.st-andrews.ac.uk/Biographies/Brahmagupta/}{\emph{Brahmagupta}, MacTutor History of Mathematics, University of St Andrews}.}
Solche Regeln ermöglichen das Rechnen mit Vorzeichen, ohne
die Zahlen zuvor als Mengen von Paaren zu definieren. Sie
beantworten zunächst die Frage, wie mit den neuen Größen
gerechnet werden soll.

Eine weitergehende Grundlagenfrage lautet: Wie lässt sich ein
Zahlenbereich so erweitern, dass jede Subtraktion möglich wird
und die bisherigen Rechengesetze erhalten bleiben? Hermann
Hankel behandelt diese Frage 1867 in seiner
\emph{Theorie der complexen Zahlensysteme}. Er führt
Differenzzeichen ein und erklärt ihre Addition durch eine
Regel, die unserer komponentenweisen Addition entspricht:
\[
  (a-b)+(c-d)=(a+c)-(b+d).
\]
Sein Prinzip der Permanenz formaler Gesetze beschreibt das
Ziel, bewährte Rechengesetze bei einer Erweiterung
weiterzuführen.\footnote{Hermann Hankel,
\emph{Theorie der complexen Zahlensysteme}, Leipzig 1867,
\S\,3 und \S\,10; zugänglich in der
\href{https://wyleyr.github.io/hankel1867tr/}{Transkription des deutschen Originaltexts mit englischer Übersetzung von Richard Lawrence}.}
Die Formel zeigt eine Verbindung zu unserem Aufbau. Sie
schreibt aber für sich noch keine Konstruktion durch
Äquivalenzklassen vor.

Die vorliegende Lesefassung beantwortet diese Grundlagenfrage
mit den Mitteln der Mengenlehre. Das Paar \((a,b)\) hält
zunächst nur fest, welche beiden natürlichen Zahlen an einer
Differenz beteiligt sein sollen. Die Bedingung \(a+d=b+c\)
legt fest, wann zwei Paare dieselbe Zahl darstellen. Durch
das Zusammenfassen zu Äquivalenzklassen entsteht der neue
Zahlenbereich; anschließend werden seine Rechengesetze
bewiesen. Die Reihenfolge dieses Aufbaus richtet sich nach
den logischen Abhängigkeiten: Zuerst brauchen wir die
Objekte, dann wohldefinierte Operationen und schließlich
ihre Gesetze. Die historischen Beispiele zeigen dagegen,
dass das Rechnen mit negativen Größen schon lange vor
einer solchen mengentheoretischen Begründung möglich war.
